\documentclass[12pt]{article}
\usepackage{geometry}
\usepackage{titling}
\usepackage{fontspec}
\usepackage{amsmath,amssymb,amsfonts}
\usepackage{unicode-math}
\usepackage{longtable}
\usepackage{booktabs}
\usepackage{xcolor}
\usepackage{nicematrix}
\usepackage{graphicx}
\usepackage[ruled,linesnumbered,noend,vlined]{algorithm2e}
\usepackage[colorlinks=true, linkcolor=blue!50!black, citecolor=blue!50!black, urlcolor=blue!50!black]{hyperref}
\usepackage{setspace}
\usepackage{tikz}
\usetikzlibrary{arrows.meta,calc}
\numberwithin{equation}{section} % equations numbered section.number

\geometry{a4paper, top=2.2cm, bottom=2.2cm, left=2.5cm, right=2.5cm}
\setlength{\droptitle}{-2cm}

\definecolor{charcoal}{HTML}{36454F}

\iftrue
\setstretch{1.025} % a little more space between lines
\setmainfont{texgyrepagella}[
  Extension = .otf,
  UprightFont = *-regular,
  BoldFont = *-bold,
  ItalicFont = *-italic,
  BoldItalicFont = *-bolditalic,
  Color=charcoal,
]
\fi

\setmathfont{texgyrepagella-math.otf}
\pagecolor{brown!10}

%
\newcommand{\scaletop}[1]{\raisebox{1.5pt}{\scalebox{0.75}{$#1\top$}}}
\newcommand{\Trans}{{\mathpalette\scaletop\relax}}
%%
\newcommand{\WOne}{\begin{pmatrix}0.5 & 0.25 \\ 0.4 & -0.1 \\ -0.3 & -0.1\end{pmatrix}}
\newcommand{\bOne}{\begin{pmatrix}0 \\ 0 \\ 0\end{pmatrix}}
\newcommand{\WTwo}{\begin{pmatrix}0.5 & -0.5 & 0.3 \\ 0.2 & 0.4 & -0.1 \\ 0.1 & 0.6 & 0.5\end{pmatrix}}
\newcommand{\bTwo}{\begin{pmatrix}0 \\ 0 \\ 0\end{pmatrix}}
\newcommand{\WThree}{\begin{pmatrix}0.3 & -0.2 & 0.5 \\ -0.4 & 0.6 & 0.1 \\ 0.2 & 0.1 & -0.3\end{pmatrix}}
\newcommand{\bThree}{\begin{pmatrix}0 \\ 0 \\ 0\end{pmatrix}}
\newcommand{\xIn}{\begin{pmatrix}1 \\ 2\end{pmatrix}}
\newcommand{\xTwoIn}{\begin{pmatrix}2 \\ 1\end{pmatrix}}
\newcommand{\etaVal}{0.1}
\newcommand{\zOne}{\begin{pmatrix}1.0000 \\ 0.2000 \\ -0.5000\end{pmatrix}}
\newcommand{\aOne}{\begin{pmatrix}1.0000 \\ 0.2000 \\ 0\end{pmatrix}}
\newcommand{\zTwo}{\begin{pmatrix}0.4000 \\ 0.2800 \\ 0.2200\end{pmatrix}}
\newcommand{\aTwo}{\begin{pmatrix}0.4000 \\ 0.2800 \\ 0.2200\end{pmatrix}}
\newcommand{\zThree}{\begin{pmatrix}0.1740 \\ 0.0300 \\ 0.0420\end{pmatrix}}
\newcommand{\pOut}{\begin{pmatrix}0.3647 \\ 0.3158 \\ 0.3196\end{pmatrix}}
\newcommand{\Lzero}{1.1528}
\newcommand{\lnThree}{1.0986}
\newcommand{\pTrue}{0.3158}
\newcommand{\dThree}{\begin{pmatrix}0.3647 \\ -0.6842 \\ 0.3196\end{pmatrix}}
\newcommand{\dTwo}{\begin{pmatrix}0.4470 \\ -0.4515 \\ 0.0180\end{pmatrix}}
\newcommand{\dOne}{\begin{pmatrix}0.1350 \\ -0.3933 \\ 0\end{pmatrix}}
\newcommand{\gWOne}{\begin{pmatrix}0.1350 & 0.2700 \\ -0.3933 & -0.7866 \\ 0 & 0\end{pmatrix}}
\newcommand{\gbOne}{\begin{pmatrix}0.1350 \\ -0.3933 \\ 0\end{pmatrix}}
\newcommand{\gWTwo}{\begin{pmatrix}0.4470 & 0.0894 & 0 \\ -0.4515 & -0.0903 & 0 \\ 0.0180 & 0.0036 & 0\end{pmatrix}}
\newcommand{\gbTwo}{\begin{pmatrix}0.4470 \\ -0.4515 \\ 0.0180\end{pmatrix}}
\newcommand{\gWThree}{\begin{pmatrix}0.1459 & 0.1021 & 0.0802 \\ -0.2737 & -0.1916 & -0.1505 \\ 0.1278 & 0.0895 & 0.0703\end{pmatrix}}
\newcommand{\gbThree}{\begin{pmatrix}0.3647 \\ -0.6842 \\ 0.3196\end{pmatrix}}
\newcommand{\epsVal}{0.01}
\newcommand{\Lplus}{1.155473}
\newcommand{\Lminus}{1.150073}
\newcommand{\fdVal}{0.270010}
\newcommand{\bpVal}{0.270010}
\newcommand{\fdRounded}{0.27000}
\newcommand{\WOneNew}{\begin{pmatrix}0.4865 & 0.2230 \\ 0.4393 & -0.0213 \\ -0.3000 & -0.1000\end{pmatrix}}
\newcommand{\bOneNew}{\begin{pmatrix}-0.0135 \\ 0.0393 \\ 0\end{pmatrix}}
\newcommand{\WTwoNew}{\begin{pmatrix}0.4553 & -0.5089 & 0.3000 \\ 0.2452 & 0.4090 & -0.1000 \\ 0.0982 & 0.5996 & 0.5000\end{pmatrix}}
\newcommand{\bTwoNew}{\begin{pmatrix}-0.0447 \\ 0.0452 \\ -0.0018\end{pmatrix}}
\newcommand{\WThreeNew}{\begin{pmatrix}0.2854 & -0.2102 & 0.4920 \\ -0.3726 & 0.6192 & 0.1151 \\ 0.1872 & 0.0911 & -0.3070\end{pmatrix}}
\newcommand{\bThreeNew}{\begin{pmatrix}-0.0365 \\ 0.0684 \\ -0.0320\end{pmatrix}}
\newcommand{\pAfter}{\begin{pmatrix}0.3190 \\ 0.4077 \\ 0.2733\end{pmatrix}}
\newcommand{\pTrueAfter}{0.4077}
\newcommand{\Lone}{0.8972}
\newcommand{\dropAct}{0.2556}
\newcommand{\dropPred}{0.2767}
\newcommand{\gradSq}{2.7666}
\newcommand{\pTwoOut}{\begin{pmatrix}0.3635 \\ 0.3684 \\ 0.2681\end{pmatrix}}
\newcommand{\LTwo}{1.3165}
\newcommand{\Lbatch}{1.2346}
\newcommand{\dThreeLabelOne}{\begin{pmatrix}-0.6353 \\ 0.3158 \\ 0.3196\end{pmatrix}}
\newcommand{\LLabelOne}{1.0088}
\newcommand{\gbThreeZero}{\begin{pmatrix}0.3333 \\ -0.6667 \\ 0.3333\end{pmatrix}}
% ---- explicit arithmetic (worked_values.tex) ----
\newcommand{\arZOne}{\begin{aligned}z^{(1)}_{1} &= 0.5\cdot 1 + 0.25\cdot 2 = 1.0000 \\ z^{(1)}_{2} &= 0.4\cdot 1 + (-0.1)\cdot 2 = 0.2000 \\ z^{(1)}_{3} &= (-0.3)\cdot 1 + (-0.1)\cdot 2 = -0.5000\end{aligned}}
\newcommand{\arZTwo}{\begin{aligned}z^{(2)}_{1} &= 0.5\cdot 1 + (-0.5)\cdot 0.2 + 0.3\cdot 0 = 0.4000 \\ z^{(2)}_{2} &= 0.2\cdot 1 + 0.4\cdot 0.2 + (-0.1)\cdot 0 = 0.2800 \\ z^{(2)}_{3} &= 0.1\cdot 1 + 0.6\cdot 0.2 + 0.5\cdot 0 = 0.2200\end{aligned}}
\newcommand{\arZThree}{\begin{aligned}z^{(3)}_{1} &= 0.3\cdot 0.4 + (-0.2)\cdot 0.28 + 0.5\cdot 0.22 = 0.1740 \\ z^{(3)}_{2} &= (-0.4)\cdot 0.4 + 0.6\cdot 0.28 + 0.1\cdot 0.22 = 0.0300 \\ z^{(3)}_{3} &= 0.2\cdot 0.4 + 0.1\cdot 0.28 + (-0.3)\cdot 0.22 = 0.0420\end{aligned}}
\newcommand{\xRow}{\begin{pmatrix}1 & 2\end{pmatrix}}
\newcommand{\aOneRow}{\begin{pmatrix}1 & 0.2 & 0\end{pmatrix}}
\newcommand{\aTwoRow}{\begin{pmatrix}0.4 & 0.28 & 0.22\end{pmatrix}}
\newcommand{\arSoftmax}{\begin{aligned}e^{z^{(3)}_1} &= e^{0.1740} = 1.1901 \\ e^{z^{(3)}_2} &= e^{0.0300} = 1.0305 \\ e^{z^{(3)}_3} &= e^{0.0420} = 1.0429 \\ \textstyle\sum_j e^{z^{(3)}_j} &= 3.2634 \\ q_1 &= 1.1901/3.2634 = 0.3647 \\ q_2 &= 1.0305/3.2634 = 0.3158 \\ q_3 &= 1.0429/3.2634 = 0.3196\end{aligned}}
\newcommand{\arLoss}{L = -\ln q_2 = -\ln 0.3158 = 1.1528}
\newcommand{\arDThree}{\begin{aligned}\delta^{(3)}_1 &= q_1 - y_1 = 0.3647 - 0 = 0.3647 \\ \delta^{(3)}_2 &= q_2 - y_2 = 0.3158 - 1 = -0.6842 \\ \delta^{(3)}_3 &= q_3 - y_3 = 0.3196 - 0 = 0.3196\end{aligned}}
\newcommand{\arDTwo}{\begin{aligned}\delta^{(2)}_{1} &= 0.3\cdot 0.3647 + (-0.4)\cdot (-0.6842) + 0.2\cdot 0.3196 = 0.4470 \\ \delta^{(2)}_{2} &= (-0.2)\cdot 0.3647 + 0.6\cdot (-0.6842) + 0.1\cdot 0.3196 = -0.4515 \\ \delta^{(2)}_{3} &= 0.5\cdot 0.3647 + 0.1\cdot (-0.6842) + (-0.3)\cdot 0.3196 = 0.0180\end{aligned}}
\newcommand{\arDOneRaw}{\begin{aligned}r^{(1)}_{1} &= 0.5\cdot 0.447 + 0.2\cdot (-0.4515) + 0.1\cdot 0.018 = 0.1350 \\ r^{(1)}_{2} &= (-0.5)\cdot 0.447 + 0.4\cdot (-0.4515) + 0.6\cdot 0.018 = -0.3933 \\ r^{(1)}_{3} &= 0.3\cdot 0.447 + (-0.1)\cdot (-0.4515) + 0.5\cdot 0.018 = 0.1883\end{aligned}}
\newcommand{\dOneRaw}{\begin{pmatrix}0.1350 \\ -0.3933 \\ 0.1883\end{pmatrix}}
\newcommand{\maskOne}{\begin{pmatrix}1 \\ 1 \\ 0\end{pmatrix}}
\newcommand{\maskTwo}{\begin{pmatrix}1 \\ 1 \\ 1\end{pmatrix}}
\newcommand{\zOneNew}{\begin{pmatrix}0.9190 \\ 0.4360 \\ -0.5000\end{pmatrix}}
\newcommand{\aOneNew}{\begin{pmatrix}0.9190 \\ 0.4360 \\ 0\end{pmatrix}}
\newcommand{\zTwoNew}{\begin{pmatrix}0.1518 \\ 0.4488 \\ 0.3499\end{pmatrix}}
\newcommand{\aTwoNew}{\begin{pmatrix}0.1518 \\ 0.4488 \\ 0.3499\end{pmatrix}}
\newcommand{\zThreeNew}{\begin{pmatrix}0.0847 \\ 0.3300 \\ -0.0701\end{pmatrix}}
\newcommand{\pOneNew}{0.3190}
\newcommand{\pThreeNew}{0.2733}
\newcommand{\pOneOld}{0.3647}
\newcommand{\pThreeOld}{0.3196}
\newcommand{\zOneB}{\begin{pmatrix}1.2500 \\ 0.7000 \\ -0.7000\end{pmatrix}}
\newcommand{\aOneB}{\begin{pmatrix}1.2500 \\ 0.7000 \\ 0\end{pmatrix}}
\newcommand{\zTwoB}{\begin{pmatrix}0.2750 \\ 0.5300 \\ 0.5450\end{pmatrix}}
\newcommand{\aTwoB}{\begin{pmatrix}0.2750 \\ 0.5300 \\ 0.5450\end{pmatrix}}
\newcommand{\zThreeB}{\begin{pmatrix}0.2490 \\ 0.2625 \\ -0.0555\end{pmatrix}}

\newcommand{\relu}{\operatorname{ReLU}}
\newcommand{\softmax}{\operatorname{softmax}}
\newcommand{\R}{\mathbb{R}}
\newcommand{\vc}[1]{\symbf{#1}}
\newcommand{\grad}[1]{\nabla_{#1}L}

\def\dx{3.5}
\def\dy{1.9}

\definecolor{relucol}{HTML}{D55E00}
\definecolor{hlcol}{HTML}{0072B2}

\title{Backpropagation with Cross-Entropy Loss\\\vspace*{4pt}\small[Draft 4]}
\author{Devi Prasad M.
  \thanks{Prepared with the assistance of Claude
  (Anthropic; model \texttt{claude-sonnet-5-5}), which produced
  the figure code, and drafts of the text
  in various places, and checked the numerical results with code.
  The author conceived the material, and reviewed and edited
  the text over many iterations. The author is responsible for
  the content.}\\Manipal School of Information Sciences.}
\date{October, 2026}

\begin{document}
\maketitle

\section{The setup}
We will use the example network of Figure~\ref{fig:network}
in this exercise. This is a fully connected network
with two hidden layers of width 3. The hidden layers use
ReLU activation. The network maps an input $\vc{x}\in\R^2$
to a probability vector $\vc{q}\in\R^3$.

\begin{figure}[htbp]
  \centering
  \resizebox{\linewidth}{!}{%
    \begin{tikzpicture}[
        font=\tiny,
        neuron/.style={circle,draw=black,line width=0.8pt,fill=white,
                        minimum size=0.95cm,inner sep=0pt},
        edge/.style={draw=black,line width=0.3pt},
        wl/.style={fill=brown!10,inner sep=0.6pt,font=\tiny},
        box/.style={draw=black,line width=0.8pt,fill=white,rounded corners=2pt,inner sep=4pt},
        arr/.style={-{Stealth[length=2.2mm]},line width=0.6pt},
    ]

    % ---------- neurons ----------
    \foreach \n/\y in {1/0.5,2/-0.5}
    \node[neuron] (i\n) at (0,\y*\dy) {$x_{\n}$};
    \foreach \n/\y in {1/1,2/0,3/-1} {
    \node[neuron] (h\n) at (\dx,\y*\dy) {};
    \node[neuron] (g\n) at (2*\dx,\y*\dy) {};
    \node[neuron] (o\n) at (3*\dx,\y*\dy) {};
    }

    % ---------- weights: edges, then labels a quarter of the way along each edge ----------
    \foreach \j in {1,2} \foreach \i in {1,2,3} \draw[edge] (i\j) -- (h\i);
    \foreach \j in {1,2,3} \foreach \i in {1,2,3} {
    \draw[edge] (h\j) -- (g\i);
    \draw[edge] (g\j) -- (o\i);
    }
    % highlighted example weight: from input 2 into hidden unit 1
    %\draw[hlcol,line width=0.6pt] (i2) -- (h1);

    \foreach \j in {1,2} \foreach \i in {1,2,3}
    \node[wl] at ($(i\j)!0.24!(h\i)$) {$W_{\i\j}$};
    \foreach \j in {1,2,3} \foreach \i in {1,2,3} {
    \node[wl] at ($(h\j)!0.24!(g\i)$) {$W_{\i\j}$};
    \node[wl] at ($(g\j)!0.24!(o\i)$) {$W_{\i\j}$};
    }
    %\node[wl,text=hlcol] at ($(i2)!0.25!(h1)$) {$W_{12}$};

    % ---------- ReLU marks (orange hinge) and identity marks ----------
    \foreach \n in {1,2,3} {
    \foreach \c in {h,g}
        \draw[relucol,line width=1.2pt]
        ($(\c\n)+(-0.23,-0.14)$) -- ($(\c\n)+(0,-0.14)$) -- ($(\c\n)+(0.21,0.16)$);
    \draw[black,line width=0.9pt] ($(o\n)+(-0.2,-0.17)$) -- ($(o\n)+(0.2,0.17)$);
    \node[font=\scriptsize] at ($(h\n)+(0,0.72)$) {\n};
    \node[font=\scriptsize] at ($(g\n)+(0,0.72)$) {\n};
    \node[font=\scriptsize] at ($(o\n)+(0,0.72)$) {\n};
    }

    % ---------- layer headings ----------
    \node[align=center] at (0,3.45)      {\scriptsize\textbf{Layer 0}\\[1pt]\footnotesize input};
    \node[align=center] at (\dx,3.45)    {\scriptsize\textbf{Layer 1}\\[1pt]\footnotesize hidden, ReLU};
    \node[align=center] at (2*\dx,3.45)  {\scriptsize\textbf{Layer 2}\\[1pt]\footnotesize hidden, ReLU};
    \node[align=center] at (3*\dx,3.45)  {\scriptsize\textbf{Layer 3}\\[1pt]\footnotesize output (logits)};

    % ---------- weight-matrix names under each bundle ----------
    \node at (0.5*\dx,-3.15)  {\scriptsize$\vc{W}^{(1)}\in\R^{3\times2}$};
    \node at (1.5*\dx,-3.15)  {\scriptsize$\vc{W}^{(2)}\in\R^{3\times3}$};
    \node at (2.5*\dx,-3.15)  {\scriptsize$\vc{W}^{(3)}\in\R^{3\times3}$};

    % ---------- softmax, probabilities, loss ----------
    \node[box,minimum width=1.4cm,minimum height=4.7cm] (sm) at (3*\dx+2.3,0) {\scriptsize softmax};
    \foreach \n in {1,2,3} \draw[arr] (o\n) --
        node[above,inner sep=1.5pt,font=\scriptsize] {$z^{(3)}_{\n}$} (sm.west |- o\n);
    \draw[arr] (sm.east) -- ++(0.55,0) node[anchor=west,inner sep=2pt] (p) {$\vc{q}$};
    \draw[arr] ($(p.east)+(0.05,0)$) -- ++(0.5,0);
    \node[box,align=center,anchor=west] (loss) at ($(p.east)+(0.6,0)$)
    {\scriptsize$L=-\ln q_c$\\[1pt]\scriptsize $c$ = label\\[2pt]\scriptsize (true class)};

    % ---------- legend ----------
    \def\ly{-4.15}
    \draw[relucol,line width=1.2pt] (0,\ly) -- ++(0.2,0) -- ++(0.2,0.3);
    \node[anchor=west] at (0.55,\ly+0.1) {\footnotesize ReLU: $a=\max(z,0)$};
    \draw[black,line width=0.9pt] (6.1,\ly-0.05) -- ++(0.4,0.3);
    \node[anchor=west] at (6.65,\ly+0.1) {\footnotesize no activation};
    %\draw[hlcol,line width=1.2pt] (7.4,\ly+0.1) -- ++(0.5,0);
    %\node[anchor=west] at (8.0,\ly+0.1)
    %{\footnotesize $W^{(1)}_{12}$: from input 2 into unit 1};
    \node[anchor=west] at (0,\ly-0.7)
    {\footnotesize $W^{(l)}_{ij}$ connects unit $j$ of layer $l\!-\!1$ to unit $i$ of layer $l$;
    biases $\vc{b}^{(l)}$ are not drawn.};

    \end{tikzpicture}}
    \caption{Example network}
  \label{fig:network}
\end{figure}

\bigskip

We assume column vectors, and define $\vc{a}^{(0)}=\vc{x}$.
The last layer, Layer 3, has no activation. Its output
$\vc{z}^{(3)}$, the \emph{logits}, goes through the softmax
function:
\begin{align}
  \vc{z}^{(l)} &= \vc{W}^{(l)}\vc{a}^{(l-1)}+\vc{b}^{(l)}, & l&=1,2,3,\label{eq:z}\\
  \vc{a}^{(l)} &= \relu(\vc{z}^{(l)}), & l&=1,2,\label{eq:a}\\
  \vc{q} &= \softmax(\vc{z}^{(3)}).\label{eq:q}
\end{align}

For $K$ classes ($K=3$ here), the softmax is defined by
\begin{equation}\label{eq:softmax}
  \softmax(\vc{z})_k=\frac{e^{z_k}}{\sum_{j=1}^{K}e^{z_j}},
  \qquad k=1,\dots,K.
\end{equation}

The weight matrix $\vc{W}^{(l)}$ of layer $l$ stores the weights of each unit as
a \emph{row}: row $i$ holds the weights into unit $i$ of layer $l$, one for each
unit of layer $l-1$ (see Figure~\ref{fig:network}). Thus $z^{(l)}_i$ in
Eq.~\eqref{eq:z} is the product of row $i$ of $\vc{W}^{(l)}$ with
$\vc{a}^{(l-1)}$, and no transpose is needed in the forward pass.

For the network of Figure~\ref{fig:network}, the three weight matrices are
\begin{equation}
  \begin{gathered}
    \vc{W}^{(1)}=\begin{bmatrix}
      W_{11} & W_{12}\\
      W_{21} & W_{22}\\
      W_{31} & W_{32}
    \end{bmatrix}\in\R^{3\times2},\\[6pt]
    \vc{W}^{(2)}=\begin{bmatrix}
      W_{11} & W_{12} & W_{13}\\
      W_{21} & W_{22} & W_{23}\\
      W_{31} & W_{32} & W_{33}
    \end{bmatrix}\in\R^{3\times3},\qquad
    \vc{W}^{(3)}=\begin{bmatrix}
      W_{11} & W_{12} & W_{13}\\
      W_{21} & W_{22} & W_{23}\\
      W_{31} & W_{32} & W_{33}
    \end{bmatrix}\in\R^{3\times3},
  \end{gathered}
  \label{eq:Wmats}
\end{equation}
where, as in the figure, the layer index is left off
the entries: $W_{ij}$ in $\vc{W}^{(l)}$ stands for
$W^{(l)}_{ij}$. Row $i$ lists the weights into
unit $i$ of layer $l$, and column $j$ the weights
leaving unit $j$ of layer $l-1$. For
example, row~1 of $\vc{W}^{(2)}$ holds $W_{11}$,
$W_{12}$, $W_{13}$, the weights from units 1, 2 and 3
of Layer~1 into unit~1 of Layer~2. Each layer also has a
bias vector $\vc{b}^{(l)}\in\R^{3}$, one entry per unit.

Each example has a
\emph{label} $c$, the index of its true class, and
a \emph{target} $\vc{y}$, the one-hot vector with
$y_c=1$ and all other entries~0. The loss is
$L=-\ln q_c$.

There are 24 weights ($6+9+9$) and 9 biases ($3+3+3$),
so 33 parameters in all, collected in $\theta$.

Training adjusts $\theta$ by gradient descent with
learning rate $\eta$, $\theta\leftarrow\theta-\eta\,\nabla_{\theta}L$.
\emph{Backpropagation} is the method for computing
all 33 partial derivatives at once, in two passes:
a forward pass that stores intermediate values,
followed by a backward sweep that reuses them.

In these notes, Part~I derives the method in general,
and Part~II demonstrates an example.

\section{The loss and its gradient at the logits}
%\paragraph{Likelihood Function.}
Given a set of observed data points
$\mathcal{D} = \{d_1, d_2, \dots, d_n\}$ assumed to be
independent and identically distributed (i.i.d.), and
a model with parameters $\theta$, the \textbf{likelihood}
function $\mathcal{L}(\theta; \mathcal{D})$ represents
the joint probability of observing that data given
those parameters:
\begin{equation}\label{eq:likelihood}
    \mathcal{L}(\theta; \mathcal{D}) = P(\mathcal{D} \mid \theta)
      = \prod_{i=1}^{n} P(d_i \mid \theta)
\end{equation}

\paragraph{The Log-Likelihood Function.}
Maximising $\mathcal{L}$ directly runs into
practical problems. Differentiating a product of
$n$ probabilities with the product rule gives a
sum of $n$ terms, each a product of $n-1$ probabilities
and one derivative, which is clumsy to compute.
\[
    \frac{d}{dx} \left(f_1(x) f_2(x) \cdots f_n(x)\right) =
        f_1'(x) f_2(x) \cdots f_n(x) +
            f_1(x) f_2'(x) \cdots f_n(x) +
                \cdots + f_1(x) f_2(x) \cdots f_n'(x)
\]
Further, multiplying many small probabilities underflows
in floating-point arithmetic. The natural logarithm
circumvents both issues. Because it is a monotonically
increasing function, it preserves the location of the
extrema. And because $\ln(ab)=\ln a+\ln b$, it turns
the product into a sum:
\begin{equation}\label{eq:loglik}
  \ell(\theta; \mathcal{D})=\ln\mathcal{L}(\theta; \mathcal{D})
    =\sum_{i=1}^{n}\ln P(d_i\mid\theta).
\end{equation}
This function $\ell(\theta; \mathcal{D})$ is the
log-likelihood. Its derivative is a sum
of $n$ simple terms.

\paragraph{Loss Function.}
In classification tasks, every data point is a labeled
example $d_i = (\symbf{x}_i, c_i)$.
By treating the input features as fixed, the model
defines the probability of the observed label as
$P(c_i \mid \symbf{x}_i, \theta) = q_{c_i}$, where
the distribution $\symbf{q}$ is predicted from
$\symbf{x}_{i}$. Summing over the dataset yields
the log-likelihood $\ell(\theta;
\mathcal{D}) = \sum_{i} \ln q_{c_i}$.
Because we minimize loss rather than maximize
likelihood, each example contributes
$-\ln q_{c_{i}}$ to the final loss function,
meaning the total loss is simply the sum of
these individual negative log-likelihood values.

For a single example, the loss $L = -\ln q_c$ is the
negative log-likelihood of the observed label $c$
given the input $\symbf{x}$ under the model
(the per-example negative log-likelihood). For a
one-hot target $\symbf{y}$, it equals the cross-entropy
\[
  H(\symbf{y}, \symbf{q}) = -\sum_{k=1}^{K} y_k \ln q_k .
\]
Averaged over a dataset $\mathcal{D}$ of $n$ examples,
the loss is $-\ell(\theta; \mathcal{D})/n$, so
minimizing it maximizes the likelihood. A network with
small random weights outputs a nearly uniform $\symbf{q}$,
so its initial loss is close to $\ln K$, where $K$ is
the number of classes.

\begin{quote}
\textit{Derivation of the initial loss $\ln K$:}
\begin{enumerate}
    \item \textbf{Near-equal logits:}
        Let $\symbf{z} = \symbf{W}\symbf{h} + \symbf{b}$
        be the final-layer scores (logits) for the
        penultimate features $\symbf{h}$. With small
        random $\symbf{W}$, $\symbf{b} = \symbf{0}$,
        and features of moderate norm, the logits
        are all close to zero:
        \[
            z_k \approx 0 \quad \forall k \in \{1, \dots, K\}.
        \]
        In fact, it suffices that $z_1 \approx \dots
        \approx z_K$, since softmax is
        invariant to adding a constant to every logit.

  \item \textbf{Uniform probabilities:} If all logits
        equal a common value $a$,
        the softmax gives
        \[
            q_k = \frac{e^{z_k}}{\sum_{j=1}^{K} e^{z_j}}
                = \frac{e^{a}}{K e^{a}} = \frac{1}{K}
            \quad \forall k,
        \]
        i.e.\ the uniform distribution over the $K$ classes.

    \item \textbf{Loss:} Substituting $q_c = 1/K$ into the per-example loss for the
    true class $c$ gives
    \[
        L = -\ln q_c = -\ln\!\left(\frac{1}{K}\right) = \ln K .
    \]
\end{enumerate}

\textit{Sanity check.} For $K = 10$ the initial loss
should be about $\ln 10 \approx 2.303$. A much larger
value usually indicates a bug or an overly large
weight initialization.
\end{quote}

When loss is calculated using natural logarithm,
\emph{nat} is the conventional unit. We may say, for
instance, ``the loss is \(2.303\) nats''.

\paragraph{Gradient of Softmax Cross-Entropy with Respect to Logits.}
Writing $\vc{z}=\vc{z}^{(3)}$, the definition of the softmax gives
$q_c=\softmax(\vc{z})_c=e^{z_c}/\sum_j e^{z_j}$, so
\begin{align}
    L &= -\ln q_c\label{eq:loss}\\[2pt]
      &= -\ln \, \dfrac{e^{z_c}}{\sum\nolimits_j e^{z_j}}\notag\\[2pt]
      &= -\bigg(\ln e^{z_c} - \ln \sum\nolimits_j e^{z_j}\bigg)\notag\\[2pt]
      &= -z_c + \ln \sum\nolimits_j e^{z_j}\, . \label{eq:lse} %%\quad \blacksquare
\end{align}
where $j$ ranges over the $K$ classes (the $K=3$ output units), and $c$ is
the fixed label.

The last line of the derivation, Eq.~\eqref{eq:lse}, is how libraries
compute $L$ (a log-sum-exp), because it avoids taking
the logarithm of a tiny probability. Note that
$L\ge0$, because $\sum_j e^{z_j}\ge e^{z_c}$.


Let $k$ range over the $K$ classes (the $K=3$ output units).
Differentiating $L=-z_c+\ln\sum_j e^{z_j}$ (Eq.~\eqref{eq:lse}) with respect
to the logit $z_k$, the first term gives $-[k{=}c]=-y_k$
(here $[k{=}c]$ is 1 if $k=c$ and 0 otherwise, which is
exactly the one-hot target), and the second gives
the softmax $q_k$:
\begin{equation}\label{eq:dLdz}
    \dfrac{\partial L}{\partial z_k}
      = -[k{=}c]+\dfrac{e^{z_k}}{\sum_j e^{z_j}}
      = - y_k + q_k = q_k - y_k \, .
\end{equation}
That is, the $k$-th entry of the gradient is the $k$-th
entry of the predicted distribution $\vc{q}$ minus the
$k$-th entry of the target $\vc{y}$. In vector form,
\begin{equation}\label{eq:delta3}
  \vc{\delta}^{(3)}:=\dfrac{\partial L}{\partial\vc{z}^{(3)}}=\vc{q}-\vc{y}\,.
\end{equation}

%\subsubsection*{From $L$ To ${\partial L}/{\partial z_k}$}
\paragraph{Working Through the Gradient Derivation.}
The chain rule is required to find the derivative of the
term $\ln {\sum_j e^{z_j}}$ with respect to $z_{k}$. Let
\begin{equation}\label{eq:u-sum-def}
    u = {\sum_j e^{z_j}} = e^{z_{1}}+e^{z_{2}}+\dots +e^{z_{K}}.
\end{equation}

Applying the chain rule, we get
\begin{equation}\label{eq:delta-intermediate-1}
\begin{aligned}[b]
    \dfrac{\partial}{\partial z_{k}}\left(\ln \sum_j e^{z_j}\right) &=
        \dfrac{\partial}{\partial z_{k}} \left(\ln u \right) \\[1pt]
        &=
        \dfrac{\ln u}{\partial u} \cdot \dfrac{\partial u}{\partial z_{k}} \\[1pt]
        &=
        \dfrac{1}{u} \cdot \dfrac{\partial u}{\partial z_{k}} \\[1pt]
        &=
        \dfrac{1}{\sum _{j}e^{z_{j}}} \cdot \dfrac{\partial}{\partial z_{k}}\left(\sum_j e^{z_j}\right).
\end{aligned}
\end{equation}

We now take the derivative of every single
term in $\sum_j e^{z_j}$, with respect to the
specific target logit \(z_{k}\):
\begin{equation}\label{eq:delta-intermediate-2}
    \dfrac{\partial}{\partial z_{k}} \left(\sum_{j}e^{z_{j}}\right) =
        \dfrac{\partial}{\partial z_{k}}\left(e^{z_{1}}+e^{z_{2}}+\dots +e^{z_{K}}\right)
\end{equation}
Note that
\begin{itemize}
  \item for any term $j$ where the index is not $k$
        (such as \(e^{z_1}, e^{z_2}\), etc.), $z_{j}$ is treated
        as a constant. The derivative of a constant is 0.
  \item for the specific term where the index is $k$
        (the term \(e^{z_{k}}\)), the derivative of \(e^{z_{k}}\)
        with respect to \(z_{k}\) is simply \(e^{z_{k}}\).
\end{itemize}

Therefore, out of the entire sum, only one term survives:
\begin{equation}\label{eq:delta-sum-logit-zk}
    \dfrac{\partial}{\partial z_{k}} \left(\sum_{j}e^{z_{j}}\right) =
        0+0+\dots +e^{z_{k}}+\dots +0=e^{z_{k}}
\end{equation}

Finally, we plug these results back in the original equation
\eqref{eq:delta-intermediate-1}
%(\(\dfrac{1}{u} \cdot \dfrac{\partial u}{\partial z_k}\)):
\begin{equation}\label{eq:delta-wrt-logit}
    \dfrac{\partial}{\partial z_{k}}\left(\ln \sum_{j}e^{z_{j}}\right)
        = \left(\frac{1}{\sum_{j}e^{z_{j}}}\right)\cdot e^{z_{k}}
        = \dfrac{e^{z_{k}}}{\sum_{j}e^{z_{j}}}
\end{equation}
Notice this is same as the softmax function. Therefore, we simply
refer to it as $q_{k}$.

\paragraph{Remarks}
Two remarks on the gradient $\vc{\delta}^{(3)}=\vc{q}-\vc{y}$
at the logits:
\begin{itemize}
  \item The entries sum to zero, since
        $\sum_k q_k=\sum_k y_k=1$. This reflects the fact that
        softmax is unchanged when the same constant is
        added to every logit, so moving all logits
        together cannot change the loss.
  \item Each entry lies in $[-1,1]$: the true
        class has $\delta_c=q_c-1\in[-1,0]$, every other
        class has $\delta_k=q_k\in[0,1]$. A gradient-descent
        step on the logits would raise the true logit in
        proportion to the probability still missing,
        $1-q_c$, and lower each wrong logit in
        proportion to the probability it holds.
\iffalse
    \item The gradient does not fade when the model is confidently wrong: as
        $q_c\to0$, $\delta_c\to-1$. A squared-error loss on the softmax outputs
        can behave differently. If the softmax saturates on a wrong class,
        $q_w\approx1$, the Jacobian factors $q_w(1-q_w)$ vanish and the
        gradient almost disappears although the prediction is badly wrong,
        whereas $\vc{q}-\vc{y}$ keeps a magnitude near~1.
\fi
\end{itemize}

\section{Backpropagation through a layer}

Consider layer $l$ with input $\vc{a}^{(l-1)}$,
and suppose
$\vc{\delta}^{(l)}=\partial L/\partial\vc{z}^{(l)}$
is known. Three results follow from the chain rule.

\paragraph{Gradient of the weights and biases.}
The pre-activation $z^{(l)}_i=\sum_j W^{(l)}_{ij}a^{(l-1)}_j+b^{(l)}_i$ depends
on $W^{(l)}_{ij}$ only through the single term $W^{(l)}_{ij}a^{(l-1)}_j$, and
no other pre-activation $z^{(l)}_{i'}$ ($i'\ne i$) involves $W^{(l)}_{ij}$,
so by the chain rule only one term survives:
\begin{equation}\label{eq:dW}
  \frac{\partial L}{\partial W^{(l)}_{ij}}
  =\frac{\partial L}{\partial z^{(l)}_i}\,\frac{\partial z^{(l)}_i}{\partial W^{(l)}_{ij}}
  =\delta^{(l)}_i\,a^{(l-1)}_j .
\end{equation}
The bias $b^{(l)}_i$ likewise enters only $z^{(l)}_i$, as a plain additive term, so
$\partial z^{(l)}_i/\partial b^{(l)}_i=1$ and
\begin{equation}\label{eq:db}
  \frac{\partial L}{\partial b^{(l)}_i}
  =\frac{\partial L}{\partial z^{(l)}_i}\,\frac{\partial z^{(l)}_i}{\partial b^{(l)}_i}
  =\delta^{(l)}_i\cdot 1=\delta^{(l)}_i .
\end{equation}
A bias behaves like a weight on an input that is always~$1$.
In matrix form, $\grad{\vc{W}^{(l)}}=\vc{\delta}^{(l)}(\vc{a}^{(l-1)})^{\top}$, an
outer product, and $\grad{\vc{b}^{(l)}}=\vc{\delta}^{(l)}$.

\paragraph{Gradient passed to the layer below.}
The entry $a^{(l-1)}_j$ influences the loss through every $z^{(l)}_i$, so the
contributions add:
\begin{equation}\label{eq:da}
  \frac{\partial L}{\partial a^{(l-1)}_j}
  =\sum_i\frac{\partial L}{\partial z^{(l)}_i}\,\frac{\partial z^{(l)}_i}{\partial a^{(l-1)}_j}
  =\sum_i\delta^{(l)}_i\,W^{(l)}_{ij}
  =\Big((\vc{W}^{(l)})^{\top}\vc{\delta}^{(l)}\Big)_j .
\end{equation}

\paragraph{Through the activation.}
Since $a^{(l-1)}_j=f(z^{(l-1)}_j)$, where $f$ is the activation function
(ReLU in layers 1 and 2) and acts elementwise,
$\partial L/\partial z^{(l-1)}_j=\big(\partial L/\partial a^{(l-1)}_j\big)\,f'(z^{(l-1)}_j)$,
and in vector form
\begin{equation}\label{eq:deltarec}
  \vc{\delta}^{(l-1)}=\Big((\vc{W}^{(l)})^{\top}\vc{\delta}^{(l)}\Big)\odot f'(\vc{z}^{(l-1)}).
\end{equation}

The backward pass is therefore: start from $\vc{\delta}^{(3)}=\vc{q}-\vc{y}$
(Eq.~\eqref{eq:delta3}); at each layer form the weight gradient as an outer
product (Eq.~\eqref{eq:dW}), multiply by $(\vc{W}^{(l)})^{\top}$
(Eq.~\eqref{eq:da}), and multiply elementwise by $f'$ of the stored
pre-activation (Eq.~\eqref{eq:deltarec}). Shapes provide a check: with $n_l$ the number of units in layer $l$,
$\vc{\delta}^{(l)}\in\R^{n_l}$ and
$\vc{a}^{(l-1)}\in\R^{n_{l-1}}$ give
$\grad{\vc{W}^{(l)}}\in\R^{n_l\times n_{l-1}}$, the shape of $\vc{W}^{(l)}$.

\paragraph{What the forward pass must store.}
For every layer, the input $\vc{a}^{(l-1)}$ (for the weight gradient) and
$\vc{z}^{(l)}$ (for $f'$). The memory needed grows with depth and width, which
is one reason why training a large network needs much more memory than running
it.

\paragraph{Cost.}
Each layer costs about one matrix--vector product in the forward pass, and
about two operations of the same size in the backward pass (the outer product
for the weight gradient and the product with $\vc{W}^{\top}$). A backward
sweep therefore costs roughly twice a forward pass, whatever the number of
parameters, whereas finite differences need two forward passes per parameter.
That is what makes backpropagation practical.

\section{The ReLU}

$\relu(z)=\max(z,0)$ has $\relu'(z)=1$ for $z>0$ and $0$ for $z<0$. At $z=0$ it
has a kink; the convention here is $\relu'(0)=0$. A hidden
unit with $z\le0$ is \emph{inactive}, and:
\begin{itemize}
  \item its $\delta$ is zero, so the weights feeding it (a row of
        $\grad{\vc{W}^{(l)}}$) and its bias receive no gradient;
  \item its output is $a=0$, so the weights leaving it (a column of the next
        layer's weight gradient, $\vc{\delta}\,\vc{a}^{\top}$) receive none
        either;
  \item a unit that is inactive for \emph{every} training example is never
        updated and cannot recover; it is called \emph{dead}.
\end{itemize}

\section{Checking a gradient numerically}

A derived gradient can be tested against the definition of a derivative. The
\emph{central difference}
\begin{equation}\label{eq:cd}
  \frac{L(\theta_i+\varepsilon)-L(\theta_i-\varepsilon)}{2\varepsilon}
  =\frac{\partial L}{\partial\theta_i}+\frac{\varepsilon^{2}}{6}\,
   \frac{\partial^{3}L}{\partial\theta_i^{3}}+\cdots
\end{equation}
has an error of order $\varepsilon^{2}$ (the one-sided difference has order
$\varepsilon$), which is why it is preferred. The step $\varepsilon$ is a
compromise: too large and the truncation error dominates; too small and the
subtraction of two nearly equal losses loses digits to rounding. In double
precision a value around $10^{-5}$ or $10^{-6}$ is typical. In a hand
calculation with $\varepsilon=0.01$, the losses must be carried to six decimals.

Each parameter needs two forward passes, so checking all $N$ parameters costs
$2N$ passes against a single backward pass for backpropagation. This is why the
check is a test of the code or the algebra, never the way a network is trained.

\section{The update, epochs and batches}

For a gradient $\vc{g}=\nabla_\theta L$, the first-order expansion gives
\begin{equation}\label{eq:firstorder}
  L(\theta-\eta\vc{g})=L(\theta)-\eta\,\lVert\vc{g}\rVert^{2}+O(\eta^{2}),
\end{equation}
so for small enough $\eta$ the loss decreases unless $\vc{g}=\vc{0}$. The
$O(\eta^{2})$ term involves curvature and explains why a large step can overshoot.

\paragraph{Epochs and batches.}
The training data consist of $n$ labelled examples, with per-example losses
$L_1,\dots,L_n$. One pass over all $n$ examples is an \emph{epoch}; training
usually runs for many epochs, often reshuffling the examples before each one.
How many updates an epoch contains depends on how many examples are used for
each update:
\begin{itemize}
  \item \emph{Per-example} (stochastic) gradient descent updates $\theta$ after
        every single example, so an epoch has $n$ updates.
  \item With \emph{minibatches} of size $B$, each update uses the gradient of the
        \emph{mean} loss over the batch, so an epoch has $n/B$ updates (rounded
        up if $B$ does not divide $n$).
  \item With $B=n$ (\emph{full-batch} gradient descent), each update uses the
        mean loss over the whole data set, so an epoch has a single update.
\end{itemize}
By linearity of differentiation, the gradient of the mean loss over a batch is
the mean of the per-example gradients. Small batches give noisier but more
frequent updates; large batches give smoother gradient estimates and make
efficient use of matrix hardware.

\paragraph{A batch as a matrix.}
With $B=5$ examples, the inputs are the columns of a matrix
$\vc{X}\in\R^{2\times5}$ and the targets are the columns of a matrix
$\vc{Y}\in\R^{3\times5}$; column $b$ of $\vc{Y}$ is the one-hot target of
column $b$ of $\vc{X}$. With labels $c=(2,3,1,2,3)$,
\[
  \vc{X}=\begin{pmatrix}
    1 & 2 & 0 & -1 & \phantom{-}3\\
    2 & 1 & 1 & \phantom{-}2 & -1
  \end{pmatrix},
  \qquad
  \vc{Y}=\begin{pmatrix}
    0 & 0 & 1 & 0 & 0\\
    1 & 0 & 0 & 1 & 0\\
    0 & 1 & 0 & 0 & 1
  \end{pmatrix}.
\]
Column $b$ of $\vc{X}$ is $\vc{x}_b$ and column $b$ of $\vc{Y}$ is
$\vc{y}_b$, so the first column of each is the pair $\vc{x}=(1,2)^{\top}$,
$\vc{y}=(0,1,0)^{\top}$ of the exercise. The forward pass keeps the same
shapes with one column per example: $\vc{A}^{(0)}=\vc{X}$, and $\vc{Z}^{(l)}$,
$\vc{A}^{(l)}$, the output $\vc{Q}$ and $\vc{\Delta}^{(3)}=\vc{Q}-\vc{Y}$ each
have $5$ columns. In general, if the examples of a batch are the columns of
$\vc{A}^{(l-1)}$ and the matching $\vc{\delta}$'s the columns of
$\vc{\Delta}^{(l)}$, the weight gradient is
$\vc{\Delta}^{(l)}(\vc{A}^{(l-1)})^{\top}$ (divided by $B$ for the mean loss):
a batch is more columns, and the same products. None of this depends on the
architecture. It holds for any feed-forward network trained by backpropagation;
only the layer sizes and the activation functions differ.

\paragraph{What practice adds.}
Large-scale training typically uses large minibatches, an adaptive optimizer
such as Adam in place of the plain step, a learning-rate schedule and gradient
clipping. The gradient formulas above are unchanged; only what is done with the
gradient differs.

\section{Summary of formulas}

\begin{itemize}
  \item \textbf{Forward:}
        $\vc{z}^{(l)}=\vc{W}^{(l)}\vc{a}^{(l-1)}+\vc{b}^{(l)}$,
        $\vc{a}^{(l)}=\relu(\vc{z}^{(l)})$,
        $\vc{q}=\softmax(\vc{z}^{(3)})$,
        $L=-\ln q_c$
        (Eqs.~\eqref{eq:z}--\eqref{eq:q}, \eqref{eq:loss}).

  \item \textbf{Start of the backward pass:}
        $\vc{\delta}^{(3)}=\vc{q}-\vc{y}$
        (Eq.~\eqref{eq:delta3}).

  \item \textbf{Each layer:}
        \[
          \begin{aligned}
            \grad{\vc{W}^{(l)}}
              &= \vc{\delta}^{(l)}(\vc{a}^{(l-1)})^{\top},\\
            \grad{\vc{b}^{(l)}}
              &= \vc{\delta}^{(l)},\\
            \vc{\delta}^{(l-1)}
              &= \bigl((\vc{W}^{(l)})^{\top}\vc{\delta}^{(l)}\bigr)
                 \odot \relu'(\vc{z}^{(l-1)}).
          \end{aligned}
        \]
        See Eqs.~\eqref{eq:dW}, \eqref{eq:db}, and \eqref{eq:deltarec}.

  \item \textbf{Check:} Central difference with step $\varepsilon$
        (Eq.~\eqref{eq:cd}).

  \item \textbf{Update:}
        $\theta\leftarrow\theta-\eta\nabla_\theta L$.
\end{itemize}

\section{Differences from PyTorch}
\label{sec:pytorch}

Code written with PyTorch computes the same numbers as
our notes here, but it does not look exactly like the
formulas above. The weights of a layer are laid out as
in the notes, whereas a batch is laid out the other way
round, and the transposes in the code come from that.

Table~\ref{tab:pytorch} lists the differences.

\begin{table}[htbp]
  \centering
  \small
  \begin{tabular}{@{}>{\raggedright\arraybackslash}p{2.6cm}>{\raggedright\arraybackslash}p{5.9cm}>{\raggedright\arraybackslash}p{6.3cm}@{}}
    \toprule
    & \textbf{Our notes} & \textbf{PyTorch} \\
    \midrule
    Weights of a layer
      & $\vc{W}^{(l)}$ of shape $n_l\times n_{l-1}$; row $i$ holds the weights into unit $i$
      & \texttt{weight} of \texttt{nn.Linear}, shape \texttt{(out\_features, in\_features)}: the same layout \\[4pt]
    Bias
      & $\vc{b}^{(l)}\in\R^{n_l}$
      & \texttt{bias}, shape \texttt{(out\_features,)} \\[4pt]
    A batch
      & examples are the \emph{columns} of $\vc{A}^{(l-1)}$, shape $n_{l-1}\times B$
      & examples are the \emph{rows} of \texttt{X}, shape \texttt{(B, in\_features)} \\[4pt]
    A layer
      & $\vc{Z}^{(l)}=\vc{W}^{(l)}\vc{A}^{(l-1)}+\vc{b}^{(l)}$ (Eq.~\eqref{eq:z})
      & \texttt{Z = X @ weight.T + bias} \\[4pt]
    Weight gradient
      & $\vc{\Delta}^{(l)}(\vc{A}^{(l-1)})^{\top}$, divided by $B$ for the mean loss
      & \texttt{weight.grad}, same shape as \texttt{weight}; includes the $1/B$ of the default mean loss \\[4pt]
    Loss
      & $L=-\ln q_c$, computed as $-z_c+\ln\sum_j e^{z_j}$ (Eq.~\eqref{eq:lse})
      & \texttt{nn.CrossEntropyLoss} takes the logits $\vc{z}^{(3)}$, not $\vc{q}$, so the network does not end in a softmax layer \\[4pt]
    Label
      & $c\in\{1,2,3\}$; target $\vc{y}$ one-hot
      & class index in $\{0,1,2\}$, no one-hot vector needed \\
    \bottomrule
  \end{tabular}
  \caption{Notes versus PyTorch conventions}
  \label{tab:pytorch}
\end{table}

The transpose in \texttt{X @ weight.T} appears because examples are stored as
rows; the weight matrix itself is the same $\vc{W}^{(l)}$. When comparing code
with these notes, read \texttt{weight} as $\vc{W}^{(l)}$ and \texttt{X} as
$(\vc{A}^{(l-1)})^{\top}$. Remember too that classes are numbered from~0 in
code, so the label $c=2$ of the worked example is the index \texttt{1}.

%----------------------------------------------------------
\clearpage
\part*{Part II: Worked computation}
%----------------------------------------------------------

\section{Data}
\[
  \vc{W}^{(1)}=\WOne,\qquad \vc{W}^{(2)}=\WTwo,
\]
\[
  \vc{W}^{(3)}=\WThree,\qquad \vc{b}^{(1)}=\vc{b}^{(2)}=\vc{b}^{(3)}=\vc{0}.
\]
Input $\vc{x}=\xIn$, label $c=2$ (so the target is $\vc{y}=(0,1,0)^{\top}$), learning
rate $\eta=\etaVal$.

\section{Forward pass and loss (Tasks 1 and 2)}

Layer 1, using Eq.~\eqref{eq:z} (the biases are zero):
\begin{equation}\label{eq:wz1}
\arZOne
\end{equation}
Applying the ReLU gives $\vc{a}^{(1)}=\aOne$: unit 3 is inactive because
$z^{(1)}_3=-0.5<0$.

Layer 2:
\begin{equation}\label{eq:wz2}
\arZTwo
\end{equation}
Because $a^{(1)}_3=0$, the third column of $\vc{W}^{(2)}$ had no effect here. All
three values are positive, so $\vc{a}^{(2)}=\vc{z}^{(2)}=\aTwo$.

Layer 3 (the logits):
\begin{equation}\label{eq:wz3}
\arZThree
\end{equation}
Softmax, Eq.~\eqref{eq:softmax}:
\begin{equation}\label{eq:wsoft}
\arSoftmax
\end{equation}
Loss, Eq.~\eqref{eq:loss}:
\begin{equation}\label{eq:wloss}
\arLoss
\end{equation}
This is slightly above $\ln 3=\lnThree$ because the true class received
$q_2=\pTrue<1/3$.

\section{Gradient at the logits (Task 3)}
By Eq.~\eqref{eq:dLdz},
\begin{equation}\label{eq:wd3}
\arDThree
\end{equation}
so $\vc{\delta}^{(3)}=\dThree$, and its entries sum to zero.

\section{Backward pass (Task 4)}

\paragraph{Layer 3.}
By Eqs.~\eqref{eq:dW} and~\eqref{eq:db},
\begin{gather}
  \grad{\vc{W}^{(3)}}=\vc{\delta}^{(3)}(\vc{a}^{(2)})^{\top}
  =\dThree\,\aTwoRow=\gWThree,\label{eq:wgW3}\\
  \grad{\vc{b}^{(3)}}=\vc{\delta}^{(3)}=\gbThree .\label{eq:wgb3}
\end{gather}

\paragraph{Layer 2.}
First pass the gradient down through $(\vc{W}^{(3)})^{\top}$
(Eq.~\eqref{eq:da}):
\begin{equation}\label{eq:wd2raw}
\arDTwo
\end{equation}
All second-layer units are active, $\relu'(\vc{z}^{(2)})=\maskTwo$, so
$\vc{\delta}^{(2)}=\dTwo$. Then
\begin{equation}\label{eq:wgW2}
  \grad{\vc{W}^{(2)}}=\vc{\delta}^{(2)}(\vc{a}^{(1)})^{\top}
  =\dTwo\,\aOneRow=\gWTwo,
\end{equation}
and $\grad{\vc{b}^{(2)}}=\vc{\delta}^{(2)}$. The third column is zero because
$a^{(1)}_3=0$.

\paragraph{Layer 1.}
The product with $(\vc{W}^{(2)})^{\top}$ (Eq.~\eqref{eq:da}) gives
\begin{equation}\label{eq:wd1raw}
\arDOneRaw
\end{equation}
Now apply the mask $\relu'(\vc{z}^{(1)})=\maskOne$ (Eq.~\eqref{eq:deltarec}):
\begin{equation}\label{eq:wd1}
  \vc{\delta}^{(1)}=\dOneRaw\odot\maskOne=\dOne .
\end{equation}
The value $0.1883$ is discarded: unit 3 is inactive, so it cannot pass gradient
on. Finally
\begin{gather}
  \grad{\vc{W}^{(1)}}=\vc{\delta}^{(1)}\vc{x}^{\top}
  =\dOne\,\xRow=\gWOne,\label{eq:wgW1}\\
  \grad{\vc{b}^{(1)}}=\vc{\delta}^{(1)}=\gbOne .\label{eq:wgb1}
\end{gather}
The third row of $\grad{\vc{W}^{(1)}}$ is zero, for the reason given above.

\section{Numerical check (Task 5)}

Take $W^{(1)}_{12}=0.25$, the weight from $x_2$ into hidden unit 1, with
$\varepsilon=\epsVal$. This entry is a good test because it influences the loss
through all later layers, so it exercises the whole backward chain. Two extra
forward passes give
\begin{equation}\label{eq:wLpm}
  L(+\varepsilon)=\Lplus,\qquad L(-\varepsilon)=\Lminus,
\end{equation}
so the central difference of Eq.~\eqref{eq:cd} is
\begin{equation}\label{eq:wcd}
  \frac{L(+\varepsilon)-L(-\varepsilon)}{2\varepsilon}=\fdVal,
\end{equation}
(or $\fdRounded$ if the printed six-decimal losses are used), against
$\bpVal$ from backpropagation, the $(1,2)$ entry of
$\grad{\vc{W}^{(1)}}$ in Eq.~\eqref{eq:wgW1}. All 33 gradient entries were checked
the same way in double precision, with a largest discrepancy below $10^{-9}$.

\section{One update (Task 6)}

With $\eta=\etaVal$, the first layer becomes
\begin{equation}\label{eq:wup1W}
  \vc{W}^{(1)}\leftarrow\WOne-\etaVal\,\gWOne=\WOneNew,
\end{equation}
\begin{equation}\label{eq:wup1b}
  \vc{b}^{(1)}\leftarrow\bOne-\etaVal\,\gbOne=\bOneNew .
\end{equation}
The other layers are updated in the same way:
\begin{align}
  \vc{W}^{(2)}&\leftarrow\WTwoNew, & \vc{b}^{(2)}&\leftarrow\bTwoNew,\\
  \vc{W}^{(3)}&\leftarrow\WThreeNew, & \vc{b}^{(3)}&\leftarrow\bThreeNew .
\end{align}
Re-running the forward pass with the new parameters:
\begin{gather}
  \vc{z}^{(1)}=\zOneNew,\quad \vc{a}^{(1)}=\aOneNew,\quad \vc{z}^{(2)}=\zTwoNew,\\
  \vc{a}^{(2)}=\aTwoNew,\quad \vc{z}^{(3)}=\zThreeNew,\quad \vc{q}=\pAfter .
\end{gather}
The new loss is $L=\Lone$, a decrease of $\dropAct$. The probability of the true
class rose from $\pTrue$ to $\pTrueAfter$, while the wrong classes fell
($q_1$ from $\pOneOld$ to $\pOneNew$, $q_3$ from $\pThreeOld$ to $\pThreeNew$).
Unit 3 of the first layer is still inactive, with $z^{(1)}_3=-0.5$ unchanged:
a dead unit does not move at all.

For the first-order check, Eq.~\eqref{eq:firstorder}, the sum of squares of
all 33 gradient entries is
$\lVert\nabla L\rVert^{2}=\gradSq$, so $\eta\lVert\nabla L\rVert^{2}=\dropPred$,
against the actual decrease of $\dropAct$. The difference is the second-order
term, and it shrinks as $\eta\to0$.

\section{Beyond one example (Tasks 7 and 8)}

\paragraph{A second example.}
For $\vc{x}'=\xTwoIn$ with label $c'=3$, under the \emph{original}
parameters:
\begin{gather}
  \vc{z}^{(1)}=\zOneB,\quad \vc{a}^{(1)}=\aOneB,\quad \vc{z}^{(2)}=\zTwoB,\\
  \vc{a}^{(2)}=\aTwoB,\quad \vc{z}^{(3)}=\zThreeB,\quad \vc{q}'=\pTwoOut,
\end{gather}
so $L'=-\ln q'_3=\LTwo$, and the mean loss over the two examples is $\Lbatch$.
The gradient of the mean loss is the mean of the two per-example gradients. (A
batch gradient was not computed here.) Unit 3 of layer 1 is inactive for both
examples; had it been inactive for \emph{all} training data, it would be dead.
With $n=4$ training examples, per-example updates give 4 updates per epoch and
minibatches of size $B=2$ give 2.

\paragraph{All parameters zero.}
Every $\vc{z}$ and $\vc{a}$ is then zero and $\vc{q}$ is uniform. The weight
gradients vanish, $\grad{\vc{W}^{(3)}}=\vc{\delta}^{(3)}(\vc{a}^{(2)})^{\top}=\vc{0}$
because $\vc{a}^{(2)}=\vc{0}$, and $\vc{\delta}^{(2)}=\vc{0}$ because
$\vc{W}^{(3)}=\vc{0}$. Only the output bias has a gradient,
\begin{equation}\label{eq:wzero}
  \grad{\vc{b}^{(3)}}=\gbThreeZero ,
\end{equation}
and over later steps the hidden layers never receive one (confirmed over 50
steps): the network can learn only the class frequencies. More generally, any
initialisation that makes the units of a layer identical keeps them identical,
since they receive identical gradients; this is why weights are initialised
randomly.

\paragraph{A different label.}
If the label were 1 instead of 2, so that $\vc{y}=(1,0,0)^{\top}$, the forward
quantities $\vc{z}^{(l)}$, $\vc{a}^{(l)}$, $\vc{q}$ of
Eqs.~\eqref{eq:wz1}--\eqref{eq:wsoft} would not change, because they do not
depend on the label. The label enters only through the loss,
$L=-\ln q_1=\LLabelOne$ (Eq.~\eqref{eq:loss}), and through
$\vc{\delta}^{(3)}=\vc{q}-\vc{y}=\dThreeLabelOne$ (Eq.~\eqref{eq:delta3}).
Everything downstream of $\vc{\delta}^{(3)}$ in the backward pass would change.


\begin{table}[htbp]
  \centering
  \small
  \renewcommand{\arraystretch}{1.15}
  \begin{tabular}{@{}>{\raggedright\arraybackslash}p{2.4cm}>{\raggedright\arraybackslash}p{5.9cm}>{\raggedright\arraybackslash}p{6.4cm}@{}}
    \toprule
    \textbf{Quantity} & \textbf{Definition} & \textbf{Common-sense meaning} \\
    \midrule
    Surprisal of outcome $k$ under $\vc{q}$
      & $I_{\vc{q}}(k)=-\ln q_k$
      & How surprised the model is when outcome $k$ occurs. It is $0$ when $q_k=1$ and grows without bound as $q_k\to0$; surprisals of independent events add. The loss $L=-\ln q_c$ is the surprisal of the true label. \\[6pt]
    Entropy of $\vc{y}$
      & $H(\vc{y})=-\sum_k y_k\ln y_k=\mathbb{E}_{k\sim\vc{y}}[-\ln y_k]$
      & The average surprisal of outcomes drawn from $\vc{y}$, judged by $\vc{y}$ itself: the uncertainty built into $\vc{y}$. It is $0$ for a one-hot $\vc{y}$ and largest, $\ln K$, for the uniform one. \\[6pt]
    Cross-entropy
      & $H(\vc{y},\vc{q})=-\sum_k y_k\ln q_k=\mathbb{E}_{k\sim\vc{y}}[-\ln q_k]$
      & The average surprisal of the model $\vc{q}$ when outcomes actually come from $\vc{y}$. This is the training loss; for a one-hot $\vc{y}$ it equals $-\ln q_c$. \\[6pt]
    KL divergence
      & $\operatorname{KL}(\vc{y}\,\|\,\vc{q})=\sum_k y_k\ln\dfrac{y_k}{q_k}=H(\vc{y},\vc{q})-H(\vc{y})$
      & The extra surprisal caused by using $\vc{q}$ instead of the truth $\vc{y}$. It is never negative, it is $0$ only when $\vc{q}=\vc{y}$, and it is not symmetric in $\vc{y}$ and $\vc{q}$. \\[6pt]
    Perplexity
      & $\mathrm{PP}=e^{H(\vc{y},\vc{q})}$ (base $e$ with $\ln$; use $2^{H}$ if $H$ is in bits)
      & The number of equally likely options the model is effectively hesitating between. A uniform $\vc{q}$ over $K$ classes has $\mathrm{PP}=K$ and a perfect model has $\mathrm{PP}=1$. In practice it is $e^{\text{mean loss}}$. \\
    \bottomrule
  \end{tabular}
  \caption{Surprisal, entropy, cross-entropy, KL divergence and perplexity}
  \label{tab:info}
\end{table}

\newpage

\paragraph{Cross-entropy $=$ entropy $+$ KL divergence.}
For each class $k$, $-\ln q_k=-\ln y_k+\ln\dfrac{y_k}{q_k}$, because
$\ln\dfrac{y_k}{q_k}=\ln y_k-\ln q_k$. Multiplying by $y_k$ and summing over $k$ gives
\begin{equation}
  \underbrace{-\sum_k y_k\ln q_k}_{H(\vc{y},\vc{q})}
  =\underbrace{-\sum_k y_k\ln y_k}_{H(\vc{y})}
  +\underbrace{\sum_k y_k\ln\frac{y_k}{q_k}}_{\operatorname{KL}(\vc{y}\,\|\,\vc{q})}.
  \label{eq:cekl}
\end{equation}
(Terms with $y_k=0$ contribute $0$ throughout, with the convention $0\ln0=0$.)

\paragraph{KL divergence is never negative.}
Use $\ln t\le t-1$ for $t>0$, with equality only at $t=1$. Taking
$t=q_k/y_k$ for the classes with $y_k>0$,
\[
  -\operatorname{KL}(\vc{y}\,\|\,\vc{q})=\sum_{y_k>0}y_k\ln\frac{q_k}{y_k}
  \le\sum_{y_k>0}y_k\Big(\frac{q_k}{y_k}-1\Big)
  =\sum_{y_k>0}q_k-\sum_{y_k>0}y_k
  =\sum_{y_k>0}q_k-1\le0 .
\]
Hence $\operatorname{KL}\ge0$, with equality only if $q_k=y_k$ for every $k$.
By Eq.~\eqref{eq:cekl}, $H(\vc{y},\vc{q})\ge H(\vc{y})$: the loss cannot go
below the entropy of the targets, and it reaches that value exactly when
$\vc{q}=\vc{y}$. For a one-hot $\vc{y}$ the floor is $H(\vc{y})=0$.

\paragraph{Maximum entropy, and why a fresh network has loss near $\ln K$.}
Let $\vc{u}$ be the uniform distribution, $u_k=1/K$. Then
\[
  \operatorname{KL}(\vc{y}\,\|\,\vc{u})=\sum_k y_k\ln\frac{y_k}{1/K}
  =\sum_k y_k\ln y_k+\ln K\sum_k y_k=-H(\vc{y})+\ln K .
\]
Since $\operatorname{KL}\ge0$, $H(\vc{y})\le\ln K$, with equality only for a
uniform $\vc{y}$. Also, for any $\vc{y}$,
$H(\vc{y},\vc{u})=-\sum_k y_k\ln\frac1K=\ln K$: a model that outputs a uniform
$\vc{q}$ has loss $\ln K$ whatever the target.

\paragraph{Perplexity is the exponential of the mean loss.}
For labelled examples $(\vc{x}_i,c_i)$, $i=1,\dots,n$, let
$\bar L=-\frac1n\sum_{i=1}^{n}\ln q_{c_i}$ be the mean loss. Then
\[
  e^{\bar L}=\exp\Big(-\frac1n\sum_{i=1}^{n}\ln q_{c_i}\Big)
  =\Big(\prod_{i=1}^{n}q_{c_i}\Big)^{-1/n},
\]
the reciprocal of the geometric mean of the probabilities given to the true
labels. For $n=1$ it is $1/q_c$: in the worked example the perplexity was
$1/0.3158\approx3.17$ before the update and $1/0.4077\approx2.45$ after it, so
the model went from hesitating between about three classes to about two and a
half.

\iffalse
\section{What to take from the computation}
\begin{itemize}
  \item The backward pass mirrors the forward pass: the same matrices, transposed,
        with the stored forward values supplying $\vc{a}^{(l-1)}$ and the ReLU masks.
  \item Gradient reaches a weight only through active units: the zero row and
        zero column above are exactly the inactive unit.
  \item One step of size $\eta=0.1$ lowered the loss from $\Lzero$ to $\Lone$ and
        moved probability toward the true class, as the first-order estimate
        predicted.
  \item A numerical derivative agreed with backpropagation to the digits shown,
        which is the standard way to trust a hand-derived gradient.
\end{itemize}
\fi

\end{document}